% Generated draft: requires mathematical and source review.
\documentclass[11pt]{article}
\usepackage[margin=25mm]{geometry}
\usepackage{fontspec}
\setmainfont{Latin Modern Roman}
\usepackage{amsmath,amssymb,mathtools,unicode-math}
\setmathfont{Latin Modern Math}
\usepackage{enumitem,needspace,xurl,hyperref}
\setlength{\parindent}{0pt}
\setlength{\parskip}{5pt}
\setlength{\emergencystretch}{3em}
\newcommand{\N}{\mathbb{N}}
\newcommand{\Z}{\mathbb{Z}}
\newcommand{\Q}{\mathbb{Q}}
\newcommand{\R}{\mathbb{R}}
\newcommand{\C}{\mathbb{C}}
\newcommand{\F}{\mathbb{F}}
\newcommand{\A}{\mathbb{A}}
\newcommand{\PP}{\mathbb{P}}
\newcommand{\E}{\mathbb{E}}
\newcommand{\eps}{\varepsilon}
\newcommand{\dd}{\,\mathrm{d}}
\begin{document}
\section*{UnsolvedMath ID: 30004663}
\textbf{Source problem code:} catalogue entry
\subsection{Definitions and mathematical statement}
For a field extension \(L/F\) and an \(F\)-variety \(V\), write \(V(L)=\operatorname{Hom}_F(\operatorname{Spec}L,V)\) for its set of \(L\)-rational points. A projective homogeneous variety is understood relative to an acting algebraic group, and a torsor is understood relative to a group under which it is a torsor; the supplied record does not specify these groups or their hypotheses.

The supplied conjecture has the following formal shape. Let \(F\) be the intended field, let \(S\) be the intended set of points corresponding to discrete valuations of \(F\), and let \(M_x\) be the intended field extension of \(F\) associated with \(x\in S\). For every projective homogeneous variety or torsor \(V\) over \(F\) in the intended class, is it true that
\[
V(F)\ne\varnothing
\quad\Longleftrightarrow\quad
\bigl(\forall x\in S,\;V(M_x)\ne\varnothing\bigr)?
\]
Here the interpretation of each \(M_x\) as an \(F\)-field is required to interpret the displayed rational-point condition; its actual construction is not supplied.

\textbf{Unresolved qualifications.} This is a statement schema, not a fully specified conjecture. The record does not give the hypotheses on \(F\), the ambient space containing the points \(x\), the precise valuation set \(S\), or the definition of \(M_x\). In particular, it does not say whether \(M_x\) is a valuation completion or another associated field. It also leaves unspecified the acting groups, the groups defining the torsors, and any rationality, smoothness, connectedness, or other restrictions needed for the intended assertion. These data cannot be recovered uniquely from the supplied statement.
\subsection{Short English statement}
The proposed local\textendash global principle asks whether every variety or torsor in the intended class has an \(F\)-rational point exactly when it has a point over every associated field \(M_x\). The field hypotheses, local fields, valuation set, and group hypotheses are missing, so the intended conjecture requires clarification. The supplied dated curation labels it partially solved; that assessment has not been independently verified.
\subsection{Sources}
\begin{itemize}
\item Local UnsolvedMath registry, permanent ID \(30004663\), source problem code \texttt{catalogue entry}, title: ``Discrete-Valuation Local-Global Principle for Homogeneous Varieties and Torsors.'' The record attributes the problem to Oberwolfach workshop discussions, gives proposed year \(2021\), and supplies no named proposer.
\item Supplied report citation: \emph{Algebraic Groups (2022)}, Oberwolfach Reports, pp.\ 1087--1148, DOI \texttt{10.4171/owr/2021/20}. The literature-source label is ``Oberwolfach Report 20/2021.'' \url{https://doi.org/10.4171/owr/2021/20}. The differing year designations are retained as supplied.
\item Supplied literature source: \emph{Patching over Berkovich curves}. \url{https://doi.org/10.1112/S0010437X19007632}. No authors, theorem numbers, or precise supporting passages are supplied.
\item Literature assessment dated exactly \texttt{2026-08-22}: supplied dated curation, not independently verified. It labels the problem ``partially solved'' and states that the local\textendash global principle is proved for projective homogeneous spaces and torsors over rational groups, and for several strong-smoothness cases, including the quadratic-form application, while the full discrete-valuation conjecture remains open. The supplied background describes this as literature triage derived from OWR curation and requests independent verification of time-sensitive status and cited claims.
\end{itemize}
\end{document}
